% Papel punteado de 5 mm (A4) — Eric Lucero González — https://ericlucero.dev
% Licencia CC BY-NC-SA 4.0. Compilar con pdfLaTeX (dos pasadas: latexmk lo hace solo).
\documentclass[a4paper]{article}
\usepackage[T1]{fontenc}
\usepackage[spanish]{babel}
\usepackage{lmodern}
\usepackage{tikz}
\pagestyle{empty}

% ── Configuración (mm) ─────────────────────────────────────
\def\anchoPapel{210}
\def\altoPapel{297}
\def\margenMin{10}   % margen mínimo alrededor de los puntos
\def\paso{5}         % distancia entre puntos
\def\radio{0.5pt}     % tamaño del punto (0.35pt = más tenue)

\begin{document}
\begin{tikzpicture}[
    remember picture, overlay, shift={(current page.south west)},
    x=1mm, y=1mm,
    dot/.style={gray!65}
  ]
  % Intervalos enteros que caben; los puntos quedan centrados en la hoja
  \pgfmathtruncatemacro{\nx}{(\anchoPapel - 2*\margenMin) / \paso}
  \pgfmathtruncatemacro{\ny}{(\altoPapel  - 2*\margenMin) / \paso}
  \pgfmathsetmacro{\x}{(\anchoPapel - \nx*\paso) / 2}
  \pgfmathsetmacro{\y}{(\altoPapel  - \ny*\paso) / 2}

  \foreach \i in {0,...,\nx} {
    \foreach \j in {0,...,\ny} {
      \fill[dot] (\x + \i*\paso, \y + \j*\paso) circle (\radio);
    }
  }
\end{tikzpicture}
\end{document}
